% ============================================================
% 30_layout_spacing.tex — Layout-Grundlagen & zentrale Spacing-Stellschrauben
% ============================================================

\setlength{\parskip}{4.0pt plus 1pt minus 0.5pt} % Mehr Luft zwischen Absätzen
\setlength{\parindent}{0pt}

\widowpenalty=10000
\clubpenalty=10000
\brokenpenalty=4000
\predisplaypenalty=10000

\raggedcolumns
\newcounter{zsfFirstChap}\setcounter{zsfFirstChap}{1}

% ------------------------------------------------------------
% SPACING-SCALE (4 Stufen)
% ------------------------------------------------------------
\newlength{\ZSFspaceXS}\setlength{\ZSFspaceXS}{1pt}
\newlength{\ZSFspaceS} \setlength{\ZSFspaceS}{2pt}
\newlength{\ZSFspaceM} \setlength{\ZSFspaceM}{4pt}
\newlength{\ZSFspaceL} \setlength{\ZSFspaceL}{8pt}

% Abstand ober-/unterhalb von Text-Box-Bindungen und Trennlinien
\newlength{\ZSFspaceOuter}  \setlength{\ZSFspaceOuter}{3pt}
\newlength{\ZSFspaceSep}    \setlength{\ZSFspaceSep}{0pt}

% ------------------------------------------------------------
% Spalten-Layout
% ------------------------------------------------------------
\setlength{\columnsep}{\ZSFspaceL} % Konsistent mit Skala
\setlength{\multicolsep}{\ZSFspaceS}
\setlength{\emergencystretch}{2em}

% ------------------------------------------------------------
% Titelbalken & Hilfsgrössen
% ------------------------------------------------------------
\newlength{\ZSFtitleBarHeight}    \setlength{\ZSFtitleBarHeight}{9.2pt}
\newlength{\ZSFtableTitleBarHeight}\setlength{\ZSFtableTitleBarHeight}{7.8pt}
\newcommand{\ZSFtableArrayStretch}{1.04}
\newlength{\ZSFtitleboxBeforeSkip}\setlength{\ZSFtitleboxBeforeSkip}{6pt} % Mehr Luft vor defbox/tablebox/figbox/warnbox
\newlength{\ZSFtitleboxAfterSkip} \setlength{\ZSFtitleboxAfterSkip}{3pt} % Abstand NACH defbox/tablebox/figbox/warnbox (siehe ZSFtightBoxGroup)
\newlength{\ZSFsubsectionBarBeforeSkip}\setlength{\ZSFsubsectionBarBeforeSkip}{8pt}     % Sichtbarer Standard-Abstand vor SubsectionBar
\newlength{\ZSFsubsectionAfterChapterGap}\setlength{\ZSFsubsectionAfterChapterGap}{1pt} % Restabstand wenn SubsectionBar direkt nach ChapterBar steht
\newlength{\ZSFsubsectionAfterGap}\setlength{\ZSFsubsectionAfterGap}{\ZSFspaceXS}        % Restabstand wenn Box direkt nach SubsectionBar steht (Box klebt an Bar)
\newlength{\ZSFtextMinPad}        \setlength{\ZSFtextMinPad}{\ZSFspaceXS}
\newlength{\ZSFpostChapterNudge}  \setlength{\ZSFpostChapterNudge}{0pt} % Kein Zusatzabstand nach ChapterBar; verhindert gestapeltes Spacing
\newlength{\ZSFbreakThreshold}    \setlength{\ZSFbreakThreshold}{2\baselineskip} % Mindesthöhe für BreakIfNotEnoughSpace
\newlength{\ZSFbreakThresholdChapter} \setlength{\ZSFbreakThresholdChapter}{5\baselineskip} % Grössere Mindesthöhe für Hauptkapitel

% Tokens für statementbox / procedure / valuegrid.
\newlength{\ZSFstatementBarWidth}\setlength{\ZSFstatementBarWidth}{2pt}
\newlength{\ZSFstatementLeftPad} \setlength{\ZSFstatementLeftPad}{\ZSFspaceM}
\newlength{\ZSFprocStepGap}      \setlength{\ZSFprocStepGap}{\ZSFspaceXS}
\newlength{\ZSFprocItemSep}      \setlength{\ZSFprocItemSep}{\ZSFspaceXS}
\newlength{\ZSFprocLeftMargin}   \setlength{\ZSFprocLeftMargin}{1.6em}
\newlength{\ZSFprocLabelSep}     \setlength{\ZSFprocLabelSep}{0.4em}
\newlength{\ZSFfactLeftMargin}   \setlength{\ZSFfactLeftMargin}{1.2em}
\newlength{\ZSFfactLabelSep}     \setlength{\ZSFfactLabelSep}{0.4em}
\newlength{\ZSFvaluegridColSep}  \setlength{\ZSFvaluegridColSep}{4pt}
\newlength{\ZSFvaluegridRowSep}  \setlength{\ZSFvaluegridRowSep}{3pt}
\newlength{\ZSFboxArc}           \setlength{\ZSFboxArc}{\ZSFspaceXS}
\newlength{\ZSFsubtleRuleWidth}  \setlength{\ZSFsubtleRuleWidth}{0.2pt}
\newlength{\ZSFemphasisRuleWidth}\setlength{\ZSFemphasisRuleWidth}{1.5pt}
\newlength{\ZSFtableFrameWidth} \setlength{\ZSFtableFrameWidth}{\ZSFsubtleRuleWidth} % Aussenrahmen der tablebox (umfasst Titelleiste + Tabelle)

% Maximale Höhe für eingebettete Grafiken: verhindert, dass Bilder
% mit hoher Aspect-Ratio (z.B. Toroid-Portrait) übermässig Platz
% einnehmen. Kombiniert mit width=\linewidth + keepaspectratio wird
% jedes Bild auf min(Spaltenbreite, ZSFfigMaxHeight) skaliert.
% Siehe \ZSFfig in styles/60_boxes.tex.
\newlength{\ZSFfigMaxHeight}      \setlength{\ZSFfigMaxHeight}{2.6cm}

% ------------------------------------------------------------
% TEXT-BOX-BINDUNGEN
% Koppeln Fliesstext direkt an Boxen — verhindert visuellen Bruch.
% ------------------------------------------------------------
\newcommand{\ZSFRobustUnskip}{%
  \ifdim\lastskip>0pt\relax\unskip\fi
  \ifnum\lastpenalty=0\else\unpenalty\fi
  \ifdim\lastskip>0pt\relax\unskip\fi
}

% Text der zur Box DARUNTER gehört (Einleitung)
\newcommand{\textVorBox}[1]{%
  \par
  \addvspace{\ZSFspaceOuter}%
  \begingroup
  \setlength{\parskip}{0pt}%
  {\small\color{TextVorBoxColor} \noindent #1\par}%
  \endgroup
  \nopagebreak
  \vspace{-\ZSFtitleboxBeforeSkip}%
  \vspace{1.5pt}%
}

% Text der zur Box DARÜBER gehört (Ergänzung/Folgerung)
\newcommand{\textNachBox}[1]{%
  \par
  \ZSFRobustUnskip
  \vspace{-\ZSFtitleboxAfterSkip}%
  \vspace{1.5pt}%
  \begingroup
  \setlength{\parskip}{0pt}%
  \nopagebreak
  {\small\color{TextNachBoxColor} \noindent #1\par}%
  \endgroup
  \addvspace{\ZSFspaceOuter}%
}

% Titel/Text vor einer Formel (innerhalb der formulabox)
\newcommand{\textVorFormel}[1]{%
  \par
  \begingroup
  \setlength{\parskip}{0pt}%
  \noindent\raggedright\textcolor{black!60}{\textbf{#1}}\par
  \endgroup
  \vspace{\ZSFspaceXS}%
  \centering
}

% Erklärungstext nach einer Formel (innerhalb der formulabox)
\newcommand{\textNachFormel}[1]{%
  \par
  \vspace{\ZSFspaceXS}%
  \begingroup
  \setlength{\parskip}{0pt}%
  \noindent\raggedright\textcolor{black!60}{\small #1}\par
  \endgroup
  \centering
}

% ------------------------------------------------------------
% Semantische Gap-Helfer für Kapiteldateien
% ------------------------------------------------------------
\newcommand{\ZSFSectionGap}{\vspace{8pt}} % Etwas mehr Luft vor Unterkapiteln
\newcommand{\ZSFgapXS}{\vspace{\ZSFspaceXS}}
\newcommand{\ZSFgapS}{\vspace{\ZSFspaceS}}
\newcommand{\ZSFgapM}{\vspace{\ZSFspaceM}}
\newcommand{\ZSFgapL}{\vspace{\ZSFspaceL}}
\newcommand{\ZSFmanualColumnBreak}{\columnbreak} % Expliziter Spaltenumbruch an gewünschter Stelle

% ------------------------------------------------------------
% Box-Gruppe ohne Zwischenabstand
% ------------------------------------------------------------
% Für Serien gleichartiger Titelboxen, die als EIN Block gelesen werden
% sollen (z. B. die kapitelweisen Symbol-Tabellen in ch00_zeichen): die
% Titelbalken übernehmen die Gruppierung, ein zusätzlicher Abstand
% zwischen den Boxen zerfasert den Block nur.
%
% Die Abstände werden lokal auf 0 gesetzt; da eine Umgebung eine
% TeX-Gruppe bildet, stellt \end{ZSFtightBoxGroup} die Defaults
% automatisch wieder her — es gibt keinen manuellen Reset, der
% vergessen werden könnte.
% Der Block soll als EINE Einheit gelesen werden: die Rahmenfarbe wird
% deshalb auf die beim Gruppenstart aktive Kapitelfarbe eingefroren, auch
% wenn einzelne Boxen darin die Palette umschalten (siehe
% \ZSFPinTableFrameColor in styles/40_colors_structure.tex).
% Aus demselben Grund werden die Titel linksbündig gesetzt: der Block wird
% wie eine Liste von oben nach unten gescannt, nicht Box für Box gelesen —
% linksbündige Titel bilden eine gemeinsame Lesekante, und rechts bleibt
% Platz für die Seitenzahl (\ZSFchapterTableTitle in styles/60_boxes.tex).
\NewDocumentEnvironment{ZSFtightBoxGroup}{}{%
  \setlength{\ZSFtitleboxBeforeSkip}{0pt}%
  \setlength{\ZSFtitleboxAfterSkip}{0pt}%
  \setlength{\ZSFsubsectionAfterGap}{0pt}%
  \ZSFPinTableFrameColor
  \tcbset{%
    zsftabletitlealign/.style={halign title=flush left},%
    zsftabletitlepadding/.style={toptitle=0.7pt, bottomtitle=0.7pt},%
    zsftabletitlecolors/.style={colbacktitle=\chaptercolor, coltitle=\chaptercolortext}%
  }%
}{}
\newcommand{\ZSFtabColGap}{\hspace{1.5em}}
